\documentclass[11pt]{tex/lawoutline}
\usepackage{amssymb}
\usepackage{enumitem}

\course{Civil Procedure}
\student{Prof: J. Greiner}
\semester{Fall 2026}

\begin{document}

\maketitle
\tableofcontents

\chapter{Exam Roadmap: Getting into Court}

\rulebox{The order of analysis}{
When a lawsuit is filed, keep the doctrines separate and ask the questions in this order:
\begin{enumerate}
    \item \textbf{Subject-matter jurisdiction (SMJ):} Can this type of case be heard in federal court?
    \item \textbf{Personal jurisdiction (PJ):} Can this court exercise power over this defendant?
    \item \textbf{Choice of law:} Which state's law will the court apply?
    \item \textbf{Removal, remand, abstention, and venue:} Can the case move, or should the court decline to proceed?
\end{enumerate}
State courts are generally courts of general jurisdiction. Federal courts need a specific statutory basis for SMJ, and a federal court still needs PJ over each defendant.
}

\examtip{Do not confuse SMJ and PJ. SMJ concerns the court's power over the kind of dispute; PJ concerns the court's power over the defendant. A federal court generally needs both.}

\examtip{On the exam, argue both sides when the doctrine is uncertain or the facts support competing inferences. For supplemental jurisdiction, use the statutory phrase ``same case or controversy'' first; the ``common nucleus of operative fact'' test is the usual way to show that relationship.}

\section{Quick Court-Access Checklist}

\rulebox{Checklist for Jurisdiction}{
The same dispute may be filed in state court even when it cannot be filed in federal court. If the case is in federal court, ask separately:
\begin{enumerate}
    \item Is there subject-matter jurisdiction?
    \item Does the court have personal jurisdiction over each defendant?
    \item Is venue proper, and is there a reason to move or decline the case?
    \item Which state's law applies?
\end{enumerate}
Satisfying one requirement does not substitute for satisfying the others.
}

\section{State and Federal Courts}

\rulebox{Basic allocation}{
State courts may hear almost every kind of claim. Federal courts may hear only cases that fall within a federal statute granting jurisdiction, such as 28 U.S.C. \S\,1331 or \S\,1332. Unless Congress makes federal jurisdiction exclusive, state and federal courts usually have concurrent jurisdiction over federal claims.
}

\note{Examples of commonly exclusive federal matters}
{Patents, copyrights, bankruptcy, and certain actions involving the United States. Always check the governing statute before assuming jurisdiction is concurrent.}

\chapter{Subject-Matter Jurisdiction}

\section{Arising-Under Jurisdiction}

\statutebox{28 U.S.C. \S\,1331 (1875)}{Federal district courts have original jurisdiction over civil actions arising under the Constitution, laws, or treaties of the United States. The statute is narrower than the broad constitutional grant in Article III.}

\note{Constitutional Boundaries}{Article III supplies the outer constitutional boundary; \S\,1331 is Congress's statutory grant within that boundary. \textit{Osborn v. Bank of the United States} reflects the broad constitutional understanding of arising-under jurisdiction, while \textit{Mottley} supplies the narrower statutory rule used in ordinary cases.}

\subsection{The Well-Pleaded Complaint Rule}

\rulebox{Well-pleaded complaint rule}{
Federal-question jurisdiction usually exists only when the plaintiff's own properly pleaded complaint establishes a federal cause of action. A federal defense, counterclaim, or anticipated defense \textbf{does not} create \S\,1331 jurisdiction.
}

\casebox{Louisville \& Nashville Railroad Co. v. Mottley}{211 U.S. 149 (1908)}{
\textbf{Takeaway:} A plaintiff cannot create federal-question jurisdiction merely by anticipating a federal defense.\par
\textbf{Facts and result:} The plaintiffs' claim was based on state contract law, although they expected the railroad to raise a federal defense. The Supreme Court held that the federal issue did not appear in the plaintiffs' well-pleaded complaint, so the federal court lacked SMJ.\par
\textbf{Rule:} Look only at the plaintiff's complaint as pleaded, not at defenses or counterclaims. The court must dismiss for lack of SMJ even if no party raises the issue.
}

\subsection{Federal Issues Embedded in State-Law Claims}

\rulebox{Smith--Grable exception}{
A state-law claim may arise under federal law when it necessarily raises a federal issue that is:
\begin{enumerate}
    \item embedded in the state-law cause of action;
    \item actually disputed;
    \begin{itemize}
        \item Don't take this for granted!
    \end{itemize}
    \item substantial to the particular case and important to the federal system; and
    \begin{itemize}
        \item  Is the resolution of this issue critical to resolving the case?
        \item Also, does this one case matter to the federal government? (Think of \textit{Merrell v. Dow)}
    \end{itemize}
    \item consistent with the federal--state balance approved by Congress.
\end{enumerate}
This is a narrow exception, not a general rule that every federal issue creates federal jurisdiction.
}

\casebox{Smith v. Kansas City Title \& Trust Co.}{255 U.S. 180 (1921)}{
\textbf{Takeaway:} A state-law claim can qualify for federal-question jurisdiction when proving the claim necessarily requires resolving a substantial and disputed federal issue.\par
Shareholders brought a state-law derivative claim challenging a corporation's purchase of federal bonds. The federal constitutional validity of the bonds was an essential issue in proving the state-law claim, so the Court found arising-under jurisdiction.
}

\casebox{Grable \& Sons Metal Products, Inc. v. Darue Engineering \& Manufacturing}{545 U.S. 308 (2005)}{
\textbf{Takeaway:} The embedded-federal-issue exception is narrow and depends on both the importance of the federal issue and the effect on the federal--state division of labor.\par
Grable's state quiet-title claim depended on whether the IRS had complied with a federal tax-notice statute. The Court allowed federal jurisdiction because the federal issue was necessarily raised, actually disputed, substantial to the federal system, and capable of being resolved without opening the federal courts to a large class of ordinary state cases.\par
\textbf{Exam factors:} Consider the federal interest, federalism and the likely number of similar cases, whether the dispute is principally legal or factual, and what Congress likely intended.
}

\examtip{For the Smith--Grable exception, do not stop after identifying a federal statute. Ask whether the federal issue is necessary to the state-law claim, actually disputed, substantial to the federal system, and compatible with the federal--state division of labor. A federal court may have power but still decide that exercising it would be imprudent.}

\casebox{Moore v. Chesapeake \& Ohio Railway Co.}{186 U.S. 642 (1903)}{
\textbf{Takeaway:} A federal issue inside a state-law claim is not automatically enough. The claim must satisfy the narrow Smith--Grable framework.
}

\casebox{Merrell Dow Pharmaceuticals, Inc. v. Thompson}{478 U.S. 804 (1986)}{
\textbf{Takeaway:} A federal question jurisdiction does not exist for a state-law claim just because it incorporates a federal standard, if Congress did not create a private federal right of action for violating that standard.
}

\subsection{Discretionary Factors}

\rulebox{Factors to look at:}
{
\begin{itemize}
    \item Is there a federal interest at stake?
    \item Balance of federal court vs state court in amount of cases
        \begin{itemize}
            \item Key sentence:  If the federal court had jurisdiction, it would sweep in a lot of other cases and overwhelm the federal court
        \end{itemize}
    \item Is this an issue of law or fact?
    \item Speculation regarding congressional intent
        \begin{itemize}
            \item Congress didn’t want it, otherwise they would have said it. 
        \end{itemize}

\end{itemize}
}

\subsection{Declaratory Judgments}

\statutebox{28 U.S.C. \S\,2201}{The Declaratory Judgment Act supplies a remedy, not an independent cause of action or an independent basis for SMJ. \par
So basically this is when someone isn't looking for damages, but wanting to court to tell someone else what to do, etc.}

\rulebox{Coercive-action or imaginary-lawsuit rule}{
To determine SMJ in a declaratory-judgment action:
\begin{enumerate}
    \item Imagine the coercive lawsuit that the opposing party would have filed if the declaratory remedy did not exist.
    \item Ask whether that imagined lawsuit would have an independent basis for federal SMJ under \S\,1331 or \S\,1332.
\end{enumerate}
If the imagined action is based only on state law, the declaratory-judgment plaintiff cannot create federal jurisdiction by filing first.
}

\casebox{First Federal Savings \& Loan Ass'n v. McReynolds}{1969}{
\textbf{Takeaway:} A declaratory-judgment action has the same jurisdictional character as the coercive action it seeks to prevent. A party cannot use a federal forum by asking for a declaration about a purely state-law dispute.
}

\subsection{Concurrent and Exclusive Jurisdiction}

\casebox{Tafflin v. Levitt}{493 U.S. 455 (1990)}{
\textbf{Takeaway:} State courts are presumed to have concurrent jurisdiction over federal claims unless Congress clearly makes federal jurisdiction exclusive.
}

\section{Diversity Jurisdiction}

\statutebox{28 U.S.C. \S\,1332}{Diversity jurisdiction requires a civil action between qualifying citizens of different states and an amount in controversy exceeding \$75,000, exclusive of interest and costs.}

\rulebox{Complete diversity}{
Under \textit{Strawbridge}, every plaintiff must be diverse from every defendant. The Constitution permits minimal diversity, but the diversity statute has been interpreted to require complete diversity.
}

\casebox{Strawbridge v. Curtiss}{7 U.S. (3 Cranch) 267 (1806)}{
\textbf{Takeaway:} For complete diversity, which is what \S\,1332 requires, no plaintiff may share state citizenship with any defendant. Check every party on both sides. 
}

\examtip{Make sure to reason through why X is domiciled in Y state, and same for corporations and unincorporated organizations.}

\subsection{Amount in Controversy}

\rulebox{Amount in controversy}{
The plaintiff's good-faith allegation controls unless it is legally impossible for the plaintiff to recover more than \$75,000. Interest and costs do not count. Claims by one plaintiff against one defendant may generally be aggregated, even if the claims are unrelated. Do not aggregate separate claims by different plaintiffs merely to reach the threshold.
}

\note{Calculations}
{When calculating the amount in controversy, consider all legally available damages, including compensatory, punitive, statutory, and injunctive relief when appropriate. If several plaintiffs seek one common and indivisible right, aggregation may be possible; separate individual injuries ordinarily cannot be combined.}

\subsection{Citizenship}

\rulebox{Natural persons: domicile}{
A natural person is a citizen of the person's domicile. Domicile requires:
\begin{enumerate}
    \item physical presence in the state; and
    \item intent to remain there indefinitely, meaning no definite plan to leave.
\end{enumerate}
A person keeps the old domicile until both residence and intent establish a new one. Temporary school, work, or travel arrangements do not necessarily change domicile.
}

\casebox{Mas v. Perry}{489 F.2d 1396 (5th Cir. 1974)}{
\textbf{Takeaway:} Domicile depends on both where a person is living and whether the person intends to remain there indefinitely. A person does not change domicile merely by moving temporarily for school or work.
}

\rulebox{Corporations}{
A corporation is a citizen of every state and foreign state where it is incorporated and of the state where it has its principal place of business. Under \textit{Hertz}, the principal place of business is the corporation's ``nerve center,'' usually where its officers direct, control, and coordinate the company's activities.
}

\rulebox{Unincorporated entities}{
Partnerships, labor unions, LLCs, and similar entities are citizens of every state of which any member is a citizen. Do not use the corporation rule for an unincorporated association.
}

\rulebox{Foreign parties}{
An alien or foreign national is treated as a citizen of the foreign state (52nd state!) for diversity purposes, subject to the statutory limits on suits involving foreign parties. 
}

\rulebox{Federally chartered banks}{A federally chartered national bank is a citizen of the state designated as its main office in its articles of association; do not automatically treat every branch as a separate citizenship.}

\rulebox{Estates and representatives}{
The legal representative of a dead person's estate has the dead person's citizenship. The representative of a minor or incompetent has the citizenship of the represented person.
}

\note{Things not covered by diversity}{Diversity does not ordinarily cover divorce, alimony, child-custody, or probate matters even when the citizenship and amount requirements appear satisfied. Check for the domestic-relations and probate exceptions before proceeding.}

\examtip{Diversity is measured when the complaint is filed. Analyze each party's citizenship, then test complete diversity, then test the amount in controversy.}

\section{Supplemental Jurisdiction}

\rulebox{Supplemental-jurisdiction sequence}{
For each additional claim, ask:
\begin{enumerate}
    \item Is there an anchoring claim with original federal jurisdiction?
    \item Does the anchoring claim have Arising Under Jursidction (AUJ)? 
    \begin{itemize}
        \item Use \S\ 1367(a): Do the claims form the same Article III case or controversy, usually because they come under the same case of controversy (share a common nucleus of operative fact)?
    \end{itemize}
    \item If the anchor is diversity-only, 
    \begin{itemize}
        \item Use \S\,1367(b): Is there complete diversity or the amount-in-controversy requirement?
    \end{itemize}
    \item Even if supplemental jurisdiction exists, should the court decline it under \S\,1367(c)?
    \begin{itemize}
        \item If the claim raises novel or complex issue of state law
        \item Added claim predominates over the anchoring claim 
        \item If the district court dismissed all claims over which it has original jurisdiction (doesn’t come up much 
        \item Other exceptional circumstances (also doesn’t come up much)
    \end{itemize}
\end{enumerate}
}

\statutebox{28 U.S.C. \S\,1367(a)}{When a federal court has original jurisdiction over at least one claim, it generally has supplemental jurisdiction over other claims that form part of the same Article III case or controversy. The usual way to show this is that the claims come under the same case of controversy (share a common nucleus of operative fact). Supplemental jurisdiction can include claims involving additional parties.}

\statutebox{28 U.S.C. \S\,1367(b)}{In a case founded solely on diversity jurisdiction, supplemental jurisdiction does not allow certain claims by plaintiffs against parties joined under Rules 14, 19, 20, or 24 when those claims would violate the requirements of \S\,1332 (Diversity + amount in controversy)

\begin{itemize}
    \item \textbf{Rule 14:} Impleader: bringing in a third party who may be liable to the defendant.
    \item \textbf{Rule 19:} Required joinder: adding a party whose participation is necessary for a complete resolution.
    \item \textbf{Rule 20:} Permissive joinder: joining additional plaintiffs or defendants with related claims.
    \item \textbf{Rule 24:} Intervention: allowing a nonparty to join an existing lawsuit.
\end{itemize}

The plaintiff cannot use supplemental jurisdiction to avoid complete diversity or the amount-in-controversy requirement. Claims by defendants, such as counterclaims or impleader claims, are generally not barred by \S\,1367(b), but they must still satisfy \S\,1367(a).

}

\statutebox{28 U.S.C. \S\,1367(c)}{A court may decline supplemental jurisdiction when the additional claim raises a novel or complex state-law issue, substantially predominates, all original-jurisdiction claims have been dismissed, or exceptional circumstances provide another compelling reason.}

\casebox{United Mine Workers v. Gibbs}{383 U.S. 715 (1966)}{
\textbf{Takeaway:} A state-law claim may travel with a federal claim when both arise from the same core facts and would ordinarily be tried together. Supplemental jurisdiction is also discretionary and should serve efficiency, fairness, and federalism.
}

\casebox{Owen Equipment \& Erection Co. v. Kroger}{437 U.S. 365 (1978)}{
\textbf{Takeaway:} In a diversity case, a plaintiff cannot use supplemental jurisdiction to add a claim against a nondiverse defendant when there is no independent basis for jurisdiction over that claim. Today, analyze this result through \S\,1367(b).
}

\casebox{Aldinger v. Howard; Finley v. United States}{427 U.S. 1 (1976); 490 U.S. 545 (1989)}{
\textbf{Takeaway:} These cases limited supplemental jurisdiction over additional parties before Congress enacted \S\,1367. The present exam analysis begins with the statute, but the cases explain why \S\,1367(a) broadly permits related claims and why \S\,1367(b) preserves limits in diversity cases.
}

\casebox{Guaranteed Systems, Inc. v. American National Can Co.}{842 F. Supp. 855 (M.D.N.C. 1994)}{
\textbf{Takeaway:} Even when an impleader claim is factually related, a diversity-only case cannot use supplemental jurisdiction to permit a plaintiff's claim against a nondiverse third-party defendant \textbf{when \S\,1367(b) makes that claim inconsistent with \S\,1332.}
}

\casebox{Exxon Mobil Corp. v. Allapattah Services, Inc.}{545 U.S. 546 (2005)}{
\textbf{Takeaway:} When at least one plaintiff satisfies complete diversity and the amount-in-controversy requirement, additional diverse plaintiffs may proceed under supplemental jurisdiction even if their individual claims are below \$75,000, so long as \S\,1367(b) does not otherwise bar them.
}

\section{Expanding the Lawsuit}

\examtip{When a problem adds a claim or party, ask two separate questions: 
\begin{enumerate}
    \item is the joinder permitted by the Federal Rules of Civil Procedure?
    \begin{itemize}
        \item FRCP 13, 14, 20, 24
    \end{itemize}
    \item is there subject-matter jurisdiction over the new claim? A claim can be procedurally permissible but still lack SMJ.
    \begin{itemize}
        \item AUJ
        \item Diversity + Amount in Controversy
        \item Supplemental J (beware 1367(b) for claims made by plaintiffs)
    \end{itemize}
\end{enumerate}
}

\section{Removal and Remand}

\statutebox{28 U.S.C. \S\,1441 (1948)}{A defendant may remove a civil action from state court to the federal district court embracing the place where the state action is pending if the federal court would have had original jurisdiction. Plaintiffs cannot remove the action they chose to file.}

\rulebox{Removal checklist}{
\begin{itemize}
    \item Only a defendant may remove.
    \item The case must be one the federal court could have heard originally.
    \item In a diversity-only case, the forum-defendant rule bars removal when a properly joined and served defendant is a citizen of the forum state.
    \item All defendants who have been properly joined and served generally must consent to removal.
    \item A notice of removal is ordinarily due within 30 days after service or receipt of the removable paper.
    \item A diversity case generally cannot be removed more than one year after commencement, subject to statutory exceptions.
\end{itemize}
}

\note{Immediacy}{Removal transfers the case to federal court immediately upon filing the notice, even if the removal is later shown to be defective. The plaintiff may move to remand for a procedural defect within 30 days; lack of SMJ can be raised at any time before final judgment, and the federal court may remand on its own.}

\casebox{Shamrock Oil \& Gas Corp. v. Sheets}{313 U.S. 100 (1941)}{
\textbf{Takeaway:} Removal is construed narrowly. A plaintiff who filed in state court cannot later remove merely because the defendant asserted a federal counterclaim.
}

\statutebox{28 U.S.C. \S\,1447}{Procedural objections to removal must generally be raised within 30 days after removal. A challenge based on lack of SMJ may be raised at any time before final judgment, and the court may notice the defect on its own.}

\note{Severance}{If a case contains a removable federal claim and a claim outside federal and supplemental jurisdiction, the federal court may sever and remand the unrelated claim rather than remand the entire action.}

\section{Abstention}

\rulebox{Colorado River abstention: basic rule}{
Federal courts generally must exercise the jurisdiction that Congress has given them. Under \textit{Colorado River}, however, a federal court may abstain in exceptional circumstances when parallel state litigation makes it substantially more efficient and appropriate for the state court to proceed.

\begin{enumerate}
    \item \textbf{Parallel proceedings:} There must be substantially the same parties and substantially the same claims in both the state and federal cases. The cases do not need to be perfectly identical.
    
    \item \textbf{Exceptional circumstances:} If parallel proceedings exist, the court has discretion to abstain. It must consider the totality of the circumstances, including:
    \begin{itemize}
        \item the purpose of any relevant congressional statute;
        \item the complexity and importance of state law compared with federal law;
        \item which court obtained jurisdiction first and how far each case has progressed;
        \item the relative convenience of the state and federal forums; and
        \item the extent to which the parties have already participated in the state litigation.
    \end{itemize}
    
    \item \textbf{Remedy:} If the dispute is better handled entirely in state court, the federal court may dismiss. If the federal court may still need to decide an issue after the state case ends, it should generally stay the federal case instead.
\end{enumerate}
}

\casebox{Colorado River Water Conservation District v. United States}{424 U.S. 800 (1976)}{
\textbf{Takeaway:} Federal courts may abstain in favor of parallel state litigation, but only in exceptional circumstances. The federal court's obligation to exercise its jurisdiction remains the starting point.

\textbf{Facts}
\begin{itemize}
    \item The United States filed a federal action seeking a declaration concerning water rights in Colorado.
    \item At the same time, a state proceeding was initiated under the McCarran Amendment to adjudicate the government's water rights along with the rights of other parties.
    \item The defendants asked the federal court to dismiss the federal action.
\end{itemize}

\textbf{Holding:} The federal court could abstain because the state and federal proceedings were substantially parallel and exceptional circumstances supported allowing the state proceeding to continue.

\textbf{Rule/Reasoning:} Courts should consider the purpose of the relevant statute, the complexity of state law, the order and progress of the two cases, the relative convenience of the forums, and the parties' participation in the state proceeding. Abstention is discretionary and should be used only when the circumstances strongly favor the state forum.

}

\casebox{Clark v. Lacy}{}{
\textbf{Takeaway:} Courts should apply the \textit{Colorado River} factors faithfully rather than turning them into a rigid, multi-part test. \par
\textbf{Practical Result:} Abstention became much easier to obtain in 7th circuit than elsewhere.
}

\chapter{Horizontal Choice of Law}

\section{The Basic HCOL Framework}

\definition{Horizontal choice of law}{The process of deciding which state's law applies when more than one state's law could govern an issue. HCOL does not decide where the case may be filed; it decides what law the chosen court will apply.}

\rulebox{HCOL exam sequence}{
\begin{enumerate}
    \item Ask whether the states' laws actually conflict on the particular issue. If there is no real conflict, apply the common rule.
    \item Decide whether the issue is procedural or substantive. 
    \begin{itemize}
        \item If procedural, the forum generally applies its own law. For this course, statutes of limitations are treated as procedural for HCOL purposes.
        \item If substantive, identify whether the forum follows the First Restatement or the Second Restatement. The forum uses its own choice-of-law method.
    \end{itemize}
\end{enumerate}
}

\rulebox{Federal courts and HCOL}{
Under \textit{Klaxon}, a federal court sitting in diversity applies the choice-of-law rules of the state in which it sits. Choice of law is issue-specific: different issues in the same case may be governed by different states' laws.
}

\section{First Restatement: Vested Rights}

\rulebox{Lex loci delicti}{
For a tort, apply the law of the place where the last event necessary to create the claim occurred---normally the place of injury. The chosen state's law governs the whole issue rather than assembling separate rules from several states.
}

\casebox{Alabama Great Southern Railroad Co. v. Carroll}{11 So. 803 (Ala. 1892)}{
\textbf{Takeaway:} Under the First Restatement, the place of injury controls the tort claim. Carroll was injured in Mississippi, so Mississippi's fellow-servant rule barred recovery even though the employment relationship and alleged negligence were connected to Alabama.
}

\note{Strengths and Criticisms of the First Remandment}{Strength: predictability and administrability. Main criticism: the place of injury may be accidental and may have little connection to the parties or the policy behind the disputed rule.}

\section{Second Restatement: Interest Analysis}

\rulebox{Second Restatement approach}{
\begin{enumerate}
    \item Start with a presumption in favor of the place of injury.
    \item Then, ask whether another state has the most significant relationship to the particular issue.
    \item Consider:
    \begin{itemize}
        \item the policies and interests of the states, 
        \item the parties' expectations, 
        \item the place of conduct and injury, 
        \item the place of the relationship, 
        \item convenience, 
        \item and the goals of the law at issue.
    \end{itemize}
\end{enumerate}
Different issues may receive different states' laws.
}

\note{Important Factors of Second Remandment}{The most important Second Restatement contacts are: 
\begin{itemize}
    \item the place of injury; 
    \item the place of conduct causing the injury; 
    \item the parties' domiciles, residences, citizenship, incorporation, and places of business; and 
    \item the place where the parties' relationship is centered. 
\end{itemize}
The approach is more responsive to interests but less predictable and may produce ``Frankenlaw'' by applying different states' laws to different issues.}

\casebox{Babcock v. Jackson}{191 N.E.2d 279 (N.Y. 1963)}{
\textbf{Takeaway:} The place of injury is not automatically controlling when it is only an accidental contact. New York law applied to a New York guest injured in Ontario because the parties, car, insurance, and trip were centered in New York, while Ontario's guest statute had little connection to this dispute.
}

\casebox{Neumeier v. Kuehner}{286 N.E.2d 454 (N.Y. 1972)}{
\textbf{Takeaway:} For guest-statute conflicts, New York supplied structured rules that generally preserve the place-of-injury result, except when the driver and passenger share a common home jurisdiction, in which case that jurisdiction's law ordinarily applies.
}

\section{Subtext}

\subsection{Does theory matter?}

\analysisbox{What the competing theories reveal}{
The First Restatement values certainty and avoids ``Frankenlaw.'' The Second Restatement values state interests and the policies behind particular rules, but may produce more discretion and forum shopping. On an exam, state the forum's method, identify the issue-specific interests, and explain the tradeoff rather than simply naming the theory.
}

\subsection{Should courts pay attention to legal academics?}

\note{How much control is there?}{Restatements summarize and influence law, but they do not automatically control. The controlling question is what method the forum's courts have adopted.}

\chapter{Personal Jurisdiction}

\examtip{The lecture emphasized that this part of personal jurisdiction contains genuine uncertainty. When Supreme Court precedent does not clearly resolve the issue, identify the competing approaches and analyze the result under each one rather than treating one approach as settled law.}

\section{What PJ Does}

\definition{Personal jurisdiction}{The court's power to bind a particular defendant or property through its judgment. PJ is separate from SMJ and must be analyzed for each defendant.}

\rulebox{State-court sequence}{
For an out-of-state defendant, ask:
\begin{enumerate}
    \item Is there a state long-arm statute authorizing jurisdiction?
    \item Does the exercise of jurisdiction satisfy the Fourteenth Amendment's Due Process Clause?
    \item If the basis is specific jurisdiction, are there minimum contacts connecting the defendant, forum, and litigation, and is jurisdiction fair and reasonable?
\end{enumerate}
}

\rulebox{Federal-court sequence}{
Under Federal Rule of Civil Procedure 4(k)(1)(A), a federal district court ordinarily has the same PJ that a court of general jurisdiction in the state where it sits would have. Thus, the federal court usually uses the forum state's long-arm statute and constitutional analysis. Rule 4(k)(2) can provide jurisdiction for a federal claim against a defendant not subject to jurisdiction in any state's courts when the defendant has sufficient contacts with the United States as a whole.
}

\definition{In personam, in rem, and quasi in rem}{
\textbf{In personam:} the judgment operates against the defendant personally.\par
\textbf{In rem:} the action determines rights in property located in the forum.\par
\textbf{Quasi in rem:} the court attaches property in the forum to support a claim against an absent defendant; historically, recovery was limited to the property's value.
}

\section{The Traditional View: Territorialism}

\casebox{Pennoyer v. Neff}{95 U.S. 714 (1877)}{
\textbf{Takeaway:} Pennoyer tied PJ to territorial sovereignty. A state could exercise general jurisdiction over a nonresident who was personally served while physically in the state, but publication alone could not support a personal judgment against an absent nonresident.\par
The case also distinguished a true property action from an in personam action. Property within the state could be seized in a proceeding directed at the property, but ownership of property alone did not create personal jurisdiction for a personal judgment.
}

\casebox{Kane v. New Jersey; Hess v. Pawloski}{242 U.S. 160 (1916); 274 U.S. 352 (1927)}{
\textbf{Takeaway:} States may require nonresident motorists using their roads to appoint an in-state agent for service or may treat road use as implied consent, especially for claims arising from that conduct. These cases are transitional examples of consent and territorial jurisdiction before International Shoe.
}

\section{Challenges to the Traditional View}

\casebox{International Shoe Co. v. Washington}{326 U.S. 310 (1945)}{
\textbf{Takeaway:} Physical presence is not required. Due process permits jurisdiction when the defendant has minimum contacts with the forum and exercising jurisdiction is consistent with traditional notions of fair play and substantial justice.\par
Washington sought unemployment contributions from International Shoe, a Delaware corporation headquartered in Missouri that used Washington-based sales representatives. Because the company's continuous and systematic activities in Washington gave rise to the claim, Washington could exercise specific jurisdiction.\par
\textbf{Two categories:} General jurisdiction does not depend on the claim's connection to the forum. Specific jurisdiction does depend on that connection.
}

\rulebox{General versus specific jurisdiction}{
\textbf{General jurisdiction:} The defendant is so at home in the forum that it may be sued there on claims unrelated to the forum. For a human, the paradigm forum is domicile. For a corporation, the paradigm forums are incorporation and principal place of business.\par
\textbf{Specific jurisdiction:} The claim arises out of or relates to the defendant's purposeful contacts with the forum. The court's power is limited to claims connected to those contacts.
}

\rulebox{General-jurisdiction checklist}{
First identify statutory authorization; in federal court, ordinarily cite Rule 4(k)(1)(A) and the forum state's long-arm statute. Then ask whether the defendant is:
\begin{itemize}
    \item domiciled in the forum;
    \item a corporation incorporated or principally based in the forum;
    \item subject to jurisdiction by consent, including a forum-selection clause or waiver; or
    \item physically present and personally served in the forum, subject to the unresolved rationale in \textit{Burnham}.
\end{itemize}
For a corporation, substantial business alone usually does not make the corporation ``at home'' under \textit{Daimler}; the affiliation must be exceptional compared with the corporation's activities everywhere else.
}

\casebox{Daimler AG v. Bauman}{571 U.S. 117 (2014)}{
\textbf{Takeaway:} General jurisdiction over an out-of-state corporation is exceptional. Continuous and substantial business in the forum is not enough by itself; the corporation must be essentially at home there relative to its overall activities.\par
The usual locations are the corporation's state of incorporation and principal place of business. Do not treat the fact that a company sells products in a state as automatically making it subject to general jurisdiction there.
}

\section{Minimum Contacts: Contracts}

\rulebox{Specific-jurisdiction test}{
After confirming statutory authorization, ask whether:
\begin{enumerate}
    \item the defendant purposefully established contacts with the forum;
    \item the contacts connect the defendant, the forum, and the subject matter of the litigation;
    \item the claim arises out of or relates to those contacts; and
    \item exercising jurisdiction is fair and reasonable.
\end{enumerate}
The court should not jump to fairness factors if there is no meaningful defendant-created contact connecting the forum to the dispute.
}

\note{Purposeful Contacts}{Purposeful contacts include directly conducting business, sending employees, advertising or providing advice in the forum, owning forum property, or otherwise reaching into the state. A plaintiff's unilateral move and the mere foreseeability that a product may travel into the state do not suffice.}

\note{Fairness}{The fairness inquiry considers the defendant's burden, the forum's interest, the plaintiff's interest and convenience, the location of witnesses and evidence, and the likely governing law. World-Wide Volkswagen emphasizes that these factors do not replace the minimum-contact requirement.}

\casebox{McGee v. International Life Insurance Co.}{355 U.S. 220 (1957)}{
\textbf{Takeaway:} One deliberate, forum-directed contract can support specific jurisdiction when the claim arises from that contract. California could hear the beneficiary's claim because the insurer deliberately assumed a California insurance relationship and received premiums from there.\par
Relevant fairness considerations included the forum's interest, the plaintiff's convenience, the connection between the claim and the forum, and the burden on the defendant.
}

\casebox{Hanson v. Denckla}{357 U.S. 235 (1958)}{
\textbf{Takeaway:} The plaintiff's or another person's unilateral connection to the forum is not enough. The defendant must purposefully avail itself of the forum's benefits and protections.\par
Florida could not exercise jurisdiction over a Delaware trustee merely because the settlor later moved to Florida. The trustee had not deliberately conducted activities in Florida. A forum's interest in helping its residents does not replace the defendant-contact requirement.
}

\section{Minimum Contacts: Products and Torts}

\casebox{World-Wide Volkswagen Corp. v. Woodson}{444 U.S. 286 (1980)}{
\textbf{Takeaway:} It is not enough that a product could foreseeably travel to the forum. The defendant must have created a relationship with the forum that makes it reasonable to anticipate being sued there.\par
The Robinsons bought a car in New York and were injured in Oklahoma while moving west. The New York retailer and regional distributor had no Oklahoma business, so Oklahoma lacked specific jurisdiction over them.\par
\textbf{Fairness factors:} defendant's contacts and burden, plaintiff's convenience, the forum's interest, the location of evidence and witnesses, and choice of law matter only after the minimum-contact link is established.
}

\casebox{J. McIntyre Machinery, Ltd. v. Nicastro}{564 U.S. 873 (2011)}{
\textbf{Takeaway:} The fractured Court did not adopt a single broad stream-of-commerce rule. The safest exam answer is that a defendant must do more than place a product into national commerce; it must purposefully target or serve the specific forum, or otherwise create a forum connection.\par
The English manufacturer sold a machine through a nationwide distributor, but did not target New Jersey, advertise there, send employees there, or pay New Jersey taxes. The plurality found no New Jersey jurisdiction. Justice Breyer's concurrence emphasized that a single isolated sale without ``something more'' was insufficient. The dissent would have given greater weight to the product's regular flow into New Jersey and the state's interest in protecting its residents.\par
\textbf{Impact:} When no opinion commands a majority, lower courts often treat the narrowest common ground---Breyer's ``something more'' approach---as the operative rule.
}

\subsection{Family Law}

\casebox{Kulko v. Superior Court}{436 U.S. 84 (1978)}{
\textbf{Takeaway:} A parent's agreement to let a child move to the other parent's state did not create personal jurisdiction for a child-support modification proceeding. In a family setting, the defendant's conduct must still create a forum connection; the child's or former spouse's move is not automatically the parent's purposeful availment.\par
California's important interest in children living there did not overcome the father's limited contacts with California or the burden of litigating across the country.
}

\section{Consent and Tag Jurisdiction}

\rulebox{Consent}{
PJ may be established by consent, including a forum-selection clause or failure to timely object. A defendant generally waives a PJ objection if it is not raised in the first Rule 12 response or motion. Consent is distinct from minimum-contacts jurisdiction.
}

\casebox{Burnham v. Superior Court}{495 U.S. 604 (1990)}{
\textbf{Takeaway:} The Court unanimously upheld jurisdiction over a nonresident who was personally served while physically present in California, even though the visit was brief and unrelated to the divorce. The Court split 4--4--1 on why.\par
Justice Scalia's plurality relied heavily on the historical tradition of tag jurisdiction. Justice Brennan's concurrence treated tradition as important evidence but also emphasized the state's benefits, the defendant's voluntary presence, and the modest burden of defending in California. Justice Stevens concurred separately, warning that the rule was broader than necessary.\par
\textbf{Exam use:} State the result confidently, then acknowledge the split if the rationale matters: physical presence plus in-state service remains a powerful basis for general jurisdiction, but the opinions disagree about whether tradition alone is sufficient.
}

\examtip{For a tag-jurisdiction problem, argue both paths: (1) under Justice Scalia's approach, historical tradition may make the tag sufficient by itself; and (2) under Justice Brennan's approach, treat the tag as one contact and complete the ordinary minimum-contacts and fairness analysis.}

\casebox{Grace v. MacArthur}{170 F. Supp. 442 (E.D. Ark. 1959)}{
\textbf{Takeaway:} A person may be personally served while traveling through the forum in an airplane. The court treated airspace over Arkansas as within Arkansas's territorial limits and found no principled reason to exempt interstate travelers from service. The decision reflects territorialism more strongly than International Shoe's fairness analysis.
}

\section{In Rem and Quasi in Rem After International Shoe}

Remember, these are lawsuits against the property, not the person!

\rulebox{Quasi in rem exam checklist}{
Remember that the historical suit is against the property, not directly against the person, and recovery is limited to the property's value. Ask whether the property is actually in the forum and whether it is corporeal or incorporeal. Then present both readings of \textit{Shaffer}:
\begin{enumerate}
    \item \textbf{Narrow reading:} quasi in rem survives, but an unusually coercive statute like Delaware's may be unconstitutional.
    \begin{enumerate}
        \item Ask if property is in forum state
        \item Ask if property is corporeal or incorporeal
        \begin{itemize}
            \item Incorporeal: walk through what would happen if courts have jurisdiction
        \end{itemize}
    \end{enumerate}
    \item \textbf{Broad reading:} all assertions of jurisdiction must satisfy \textit{International Shoe}:
    \begin{itemize}
        \item Is property in forum state? 
        \item Are there minimal contacts among the defendant, property, forum, and claim?
        \item Is QIRJ fair and reasonable?
    \end{itemize}
\end{enumerate}
Real property has a stronger and more permanent forum connection than movable property or an intangible such as a debt or stock.
}

\casebox{Harris v. Balk}{198 U.S. 215 (1905)}{
\textbf{Takeaway:} Under the old territorial approach, a debt was treated as following the debtor. Maryland could garnish the debt when the garnishee was personally served there, even though the creditor, debtor, and underlying transaction were connected to North Carolina.\par
The case illustrates the difficulty of locating intangible property and the old emphasis on power over the person who held the obligation. It also shows why courts must decide where an intangible is located before deciding whether it can be attached.\par
Here, the \textbf{black letter holding} is that debt is property, and that debt moves around with the debtor, which is still \textbf{good law}. This means that bank accounts, which count as debt, are property. 
}

\casebox{Shaffer v. Heitner}{433 U.S. 186 (1977)}{
\textbf{Takeaway:} Property in the forum is not automatically enough for quasi in rem jurisdiction. All assertions of state-court jurisdiction over a person or property must satisfy the minimum-contacts framework of International Shoe.\par
Delaware seized stock owned by directors and officers of a Delaware corporation in a shareholder derivative action. The defendants' ownership of stock alone did not create the necessary connection between the defendants, the litigation, and Delaware.\par
\textbf{Competing readings:} One view treats Shaffer as invalidating Delaware's unusually coercive sequestration statute but preserving ordinary in rem and quasi in rem jurisdiction. The broader view applies minimum contacts to all in rem and quasi in rem cases. On an exam, identify which reading the facts support and explain the relationship among the defendant, property, forum, and claim.\par
\textbf{Concurrences and dissent:} Justice Powell agreed that the Delaware statute was unconstitutional but left room for jurisdiction when property is closely tied to the dispute. Justice Stevens emphasized the practical significance of owning stock in a Delaware corporation. Justice Brennan's dissent argued that the defendants' roles as directors of a Delaware corporation created a sufficient Delaware relationship and that the majority overstated the lack of contacts.
}

\examtip{For a quasi in rem problem, give both analyses because Shaffer's broad language has not been fully clarified: (1) ordinary quasi in rem may remain available if the state's procedure is not unusually coercive; or (2) all jurisdictional assertions must satisfy International Shoe, requiring contacts among the defendant, the property, the forum, and the claim, plus fairness. Then explain which approach better serves the client.}

\section{Internet Contacts and Intentional Torts}

\casebox{Griffis v. Luban}{646 N.W.2d 526 (Minn. 2002)}{
\textbf{Takeaway:} Knowledge that an online statement will injure someone who lives in the forum is not enough by itself. For an intentional tort, the plaintiff must identify specific conduct showing that the defendant expressly aimed the conduct at the forum.\par
Luban's online statements intentionally targeted Griffis, whom she knew lived in Alabama, but the record did not show that Luban targeted Alabama itself. The Alabama judgment therefore could not be enforced in Minnesota because Alabama lacked PJ.
}

\note{Connection with the forum}{Walden v. Fiore later reinforced the forum-focus requirement: the plaintiff cannot be the only link between the defendant and the forum. Ask whether the defendant's conduct itself created a connection with the forum, not merely whether the plaintiff felt the injury there.}

\examtip{For an Internet intentional-tort question, do not stop at the fact that the defendant knew the plaintiff lived in the forum or would be injured there. Look for conduct expressly aimed at the forum itself. Fr example, conduct tied to the forum's activities, audience, market, or a relationship that is uniquely connected to that state. (Think of the slander against that one Hollywood actress saying she was a drunk)}

\section{Personal Jurisdiction Checklists}

\rulebox{In Personam Jurisdiction: exam checklist}{
\begin{enumerate}
    \item \textbf{Find statutory authorization.}
    \begin{itemize}
        \item A state needs a long-arm statute before it can exercise jurisdiction over an out-of-state defendant.
        \item In federal court, apply the long-arm statute of the state where the federal court sits under Rule 4(k)(1)(A).
        \item If the statute is more specific than the Constitution, apply its text carefully; that specificity is likely an issue spotter.
        \item If the statute extends to the constitutional limit, say so and proceed to the Due Process analysis. A general long-arm statute grants jurisdiction only up to the constitutional limit; a specific statute identifies an additional state interest in particular categories of disputes.
    \end{itemize}

    \item \textbf{Ask whether the court has general or specific jurisdiction.}

    \item \textbf{General jurisdiction:} Ask whether the defendant is ``at home'' in the forum.
    \begin{itemize}
        \item \textbf{Person:} Is the defendant domiciled in the state? Ask whether the person established a new residence and intended to remain there indefinitely.
        \item \textbf{Corporation:} Under \textit{Daimler}, the usual ``at home'' forums are the state of incorporation and principal place of business. Extensive sales or business in the forum are not automatically enough; an exceptional relationship with the forum is required.
        \item \textbf{People and \textit{Daimler}:} \textit{Daimler} concerned a corporation. If its ``at home'' reasoning is extended to people, identify the uncertainty and explain both sides.
        \item \textbf{Consent:} Has the defendant consented through a forum-selection clause, or waived the objection by failing to raise it promptly?
        \item \textbf{Tag jurisdiction:} Was the defendant personally served while physically present in the forum? Under \textit{Burnham}, discuss both the historical-tradition approach and the minimum-contacts/fairness approach.
    \end{itemize}

    \item \textbf{Specific jurisdiction:} The analysis has two separate stages.
    \begin{enumerate}
        \item \textbf{Minimum contacts and purposeful availment:} Are the defendant, the forum, and the subject matter of the litigation meaningfully connected? The contacts must be defendant-created and must arise out of or relate to the claim. Use both phrases---``minimum contacts'' and ``purposeful availment.'' Be cautious about relying only on advertisements, unilateral acts by the plaintiff, or the mere possibility that a product could enter the forum. \textit{Ford} emphasizes that the claim must sufficiently relate to the defendant's forum contacts even when it did not directly arise from the exact transaction in the forum.
        \item \textbf{Fairness and reasonableness:} Only after finding minimum contacts, consider the \textit{McGee} factors:
        \begin{itemize}
            \item the burden on the defendant;
            \item the forum state's interest in resolving the dispute;
            \item the plaintiff's interest in convenient and effective relief;
            \item the location of witnesses and evidence; and
            \item the connection between the forum, the claim, and the governing law.
        \end{itemize}
        The fairness factors cannot replace the required defendant-created connection to the forum.
    \end{enumerate}
\end{enumerate}
}

\rulebox{Quasi-in-Rem Jurisdiction: exam checklist}{
First identify that the suit is directed at the property, not directly at the person. Then discuss both possible readings of \textit{Shaffer}.

\begin{enumerate}
    \item \textbf{Traditional or narrow reading (associated with Justice Scalia):}
    \begin{itemize}
        \item Is the property located in the forum state?
        \item Is the property corporeal or incorporeal?
        \item If the property is incorporeal, explain where the law treats it as located and what the court's attachment of that intangible property would accomplish.
        \item Under this reading, quasi-in-rem jurisdiction remains available, although an unusually coercive procedure may still be unconstitutional.
    \end{itemize}

    \item \textbf{International Shoe reading (associated with Justice Brennan):}
    \begin{itemize}
        \item Is the property located in the forum state?
        \item Are there minimum contacts connecting the defendant, the property, the forum, and the claim?
        \item Is exercising jurisdiction fair and reasonable?
    \end{itemize}
\end{enumerate}

Under either approach, explain the nature of the property and the relationship among the defendant, the property, the forum, and the litigation. Historically, recovery in a quasi-in-rem action was limited to the value of the property attached.
}

\chapter{Venue}

\note{Vocabulary}
{
\begin{itemize}
    \item We say that venue is ``proper'', etc.  
\end{itemize}
}

\definition{Venue}{Venue asks which federal district is the proper geographic location for the case. Unlike SMJ and PJ, venue is primarily statutory.}

\statutebox{28 U.S.C. \S\,1391}{
Venue is proper in a federal district if:
\begin{enumerate}
    \item \S\,1391(b)(1) Venue based on residence: All defendants reside in the same state and at least one defendant resides in the district. (Only when all $\Delta$ live in the same state.) 
    \begin{itemize}
        \item \textbf{Residence}
        \item Here, a person's residence is where they are domiciled. 
        \item for a corporation (no distinction with unincorporated orgs.) they reside in any judicial district where they have personal jurisdiction $\leftarrow$ This means any personal jurisdiction.
        \item For a state in multiple districts: Choose the district where Secretary of the State is. (If incorporated in a multiple district state.)
        \item Non U.S. citizens can be sued in any judicial district. (Should sue in a place with no airport.) Joinder of the $\Delta$ can be disregarded when $\exists$ other $\Delta$. 
        \note{}{Still need PJ for everyone!}
    \end{itemize} 
    \item \S\,1391(b)(2) Venue based on events: A substantial part of the events or omissions giving rise to the claim occurred in the district, or a substantial part of the property at issue is located there; or
    \item \S\,1391(b)(3) if neither of those options works, any defendant is subject to personal jurisdiction in the district. (Hardly ever use this. Very rarely on the exam.)
\end{enumerate}
The third option is a fallback and usually matters only when the events occurred outside the United States. \par
But remember that to hear the merits of a case, courts must be able to lay venue and have personal jurisdiction over all defendants.

\note{Improper venue outside of this statute}{Nothing outside of this statute makes venue improper. If you comply with this statute, venue is okay.}
}

\rulebox{Venue checklist}{
\begin{enumerate}
    \item Identify every defendant's residence for venue purposes. A natural person resides where domiciled; a corporation or unincorporated entity resides in any district where it is subject to personal jurisdiction for the action.
    \item If all defendants reside in the same state, ask whether at least one resides in the district.
    \item If not, ask where a substantial part of the events, omissions, or property is located.
    \item If neither provision applies, ask whether any defendant is subject to personal jurisdiction in the district.
\end{enumerate}
}

\note{Removed cases}{
\textbf{Black Letter Law:} For purposes of this course, assume that venue is proper whenever a case is removed from state court to federal court. State-law venue will not be tested here. \par
This does not necessarily mean that they have PJ. (Still need to check this.)
}

\statutebox{28 U.S.C. \S\S\,1404(a), 1406(a), and 1631}{
These provisions are collectively called the \textbf{transfer statutes}. They govern whether a federal case should be transferred to another federal district or dismissed. For this course, refer to them collectively whenever a case needs to be transferred.

\textbf{Statutory orientation:} Section 1404(a) generally addresses transfer from a proper venue to another proper venue; \S\,1406(a) addresses an improper venue; and \S\,1631 addresses a want of jurisdiction.
}

\definition{Happy court}{A federal court that has personal jurisdiction over the defendants \textbf{and} proper venue. Only a happy court can hear the merits of the case. Must transfer or dismiss if the court is unhappy.}

\section{Forum Non Conveniens}

\doctrine{Forum non conveniens}
{
A court may resist imposition upon its jurisdiction even when jurisdiction is authorized by the letter of a general venue statute. The doctrine allows a court to dismiss when another adequate forum would be substantially more convenient. \textit{Gulf Oil Corp. v. Gilbert}.

\note{}
{
When forum non conveniens works, the case is dismissed, not transferred. So, we use this when we can't transfer the case. \par
Works when:
\begin{itemize}
    \item Federal court to foreign court
    \item Federal court to state court due to forum selection clause
    \item \todoitem{Other cases}
\end{itemize}
}
}

\commonlaw{Forum non conveniens Factors (Federal Common Law)}
{
First, apply presumption of deferral to $\pi$'s choice of venue. Check if there is an adequate alternative forum. Then, consider the factors that might overcome the presumption. 

Private- and public-interest factors:

\begin{enumerate}
    \item \textbf{Private-interest factors:} These concern the convenience of the parties and the practical ease of conducting a fair trial.
    \begin{itemize}
        \item Relative ease of access to sources of proof;
        \item Availability of compulsory process for unwilling witnesses;
        \item Cost of obtaining the attendance of willing witnesses;
        \item Possibility of viewing the premises, if appropriate;
        \item Other practical problems affecting whether the trial will be easy, expeditious, and inexpensive;
        \item Enforceability of any judgment obtained; and
        \item Whether the plaintiff selected an inconvenient forum to vex, harass, or oppress the defendant.
    \end{itemize}

    \item \textbf{Public-interest factors:} These concern the relationship between the forum, the dispute, and the administration of justice.
    \begin{itemize}
        \item Administrative difficulties caused by congested courts;
        \item Avoiding the burden of jury duty on a community with little connection to the litigation;
        \item The local interest in having localized controversies decided at home;
        \item The interest of affected members of the community in being able to observe the trial; and
        \item The appropriateness of having a diversity case decided by a court familiar with the state law that governs the dispute.
    \end{itemize}
\end{enumerate}

\textbf{Balancing rule:} The plaintiff's choice of forum should rarely be disturbed unless the private- and public-interest factors strongly favor the defendant. The trial court has substantial discretion, and its decision is reviewed for abuse of discretion. \textit{Gulf Oil Corp. v. Gilbert}, 330 U.S. 501 (1947).
}

\casebox{\textit{Piper Aircraft Co. v. Reyno}}{454 U.S. 235 (1981)}
{
\textbf{Takeaway}
\begin{itemize}
    \item A court may dismiss for forum non conveniens when an adequate alternative forum exists and the private- and public-interest factors strongly favor that forum.
    \item (Black letter law) A foreign plaintiff's choice of a United States forum receives less deference than a domestic plaintiff's choice of a home forum.
    \item The possibility that the alternative forum will apply less favorable law ordinarily does not prevent dismissal. It matters only if the alternative forum provides no meaningful remedy at all.
    \item Private factors favored Scotland because key witnesses, evidence, potential third-party defendants, and the accident site were located there.
    \item Public factors favored Scotland because the accident and victims were Scottish, while Pennsylvania had little connection to the dispute.
\end{itemize}
}

\section{Transfers}

When moving from one federal district court to another. 

\rulebox{Motions to Transfer Questions}
{
\begin{enumerate}
    \item \textbf{Must the original court transfer or dismiss?} If the original court lacks personal jurisdiction or venue is improper, it is an ``unhappy court'' and cannot hear the merits. It must dismiss or transfer the case. See \textit{Goldlawr}. \par 
    You must bring it up quickly, otherwise you are thought to have waived it. 
    \item \textbf{May the court transfer or must it dismiss?} Transfer is available only to a district in which the action could have been brought when it was filed. (This means there is another happy court $\leftarrow$ transferee court.) If no such district exists, the original court must dismiss. See \textit{Hoffman}.
    \item \textbf{Should the court transfer?} If multiple happy courts are available, the court asks whether transfer would be in the interest of justice. Apply the forum non conveniens factors from \textit{Gulf Oil} and \textit{Piper}.
    \item \textbf{What happens after transfer?} If the transferor court was happy, its choice-of-law rules generally follow the case. If the transferor court was unhappy, the transferee court applies its own choice-of-law rules.
\end{enumerate}
}

\subsection{Must the original court transfer or dismiss?}

\casebox{\textit{Goldlawr, Inc. v. Heiman}}{369 U.S. 463 (1962)}
{
\textbf{Takeaway}
\begin{itemize}
    \item A federal court may transfer rather than dismiss when venue is improper or when it lacks personal jurisdiction over the defendants.
    \item The transfer statutes can cure defects in both venue and personal jurisdiction.
    \item If a proper transferee court exists, transfer may better serve the interest of justice than dismissal and refiling.
\end{itemize}
}

\subsection{May the court transfer or dismiss?}

\casebox{\textit{Hoffman v. Blaski}}{363 U.S. 335 (1960)}
{
\textbf{Takeaway}
\begin{itemize}
    \item A case may be transferred only to a district where the plaintiff could have brought the action when it was originally filed.
    \item A defendant's later consent to personal jurisdiction or venue cannot make an otherwise unavailable district eligible for transfer.
    \item (Black Letter Law) The transferee court must have been a legally available forum at the time of filing, not merely at the time of transfer.
\end{itemize}

\textbf{Majority versus dissent}
\begin{itemize}
    \item \textbf{Majority:} The phrase ``where it might have been brought'' refers to a district where the plaintiff had a right to sue when the case began.
    \item \textbf{Dissent:} Venue and personal jurisdiction are personal privileges of the defendant that may be waived. If the defendant consents, the court should be able to transfer the case when convenience and justice favor transfer.
\end{itemize}
}

\subsection{If the federal court transfers, what consequences are there for the applicable law?}

\subsubsection{Horizontal Choice of Law in Diversity Cases}

\begin{itemize}
    \item \textbf{Current rule:} The federal district court uses federal law for matters of procedure and state law for matters of substance.
    \item But, we need a way to separate things into either \textbf{substance} or \textbf{procedure}. $\leftarrow$ \textbf{``Vertical Choice of Law''}.
    \begin{enumerate}
        \item Dividing the line between ``substance'' and ``procedure'' for HCoL $\neq$ dividing the line for VCoL.
        \note{Statutes of Limitations}{For our purposes in class, we will say that all limitations are procedural for horizontal purposes.}
        \item A federal court sitting in diversity uses the HCoL rules of the state in which it sits to resolve conflicts of state law. This means choice-of-law rules are substantive for VCoL purposes, and state law governs.
    \end{enumerate}
    \item \textbf{Put together:} A federal court sitting in diversity will apply the choice-of-law rules of the state in which it sits.
    \note{Statute of Limitations}{This means that a federal court sitting in diversity will apply the limitations period of the state in which it sits.}
\end{itemize}

\casebox{\textit{Van Dusen v. Barrack}}{376 U.S. 612 (1964)}
{
\textbf{Takeaway}
\begin{itemize}
    \item If venue is proper in the transferor court, the transferor court had personal jurisdiction, and the defendant requested the transfer, the transferee court applies the transferor state's choice-of-law rules.
    \item The transfer should change only the physical location of the courtroom, not the applicable state law.
\end{itemize}
}

\casebox{\textit{Ferens v. John Deere Co.}}{494 U.S. 516 (1990)}
{
\textbf{Takeaway}
\begin{itemize}
    \item The \textit{Van Dusen} rule applies even when the plaintiff requests the transfer.
    \item If the transferor court was happy, the transferor state's choice-of-law rules follow the case.
    \item If the transferor court was unhappy, the transferee court applies its own choice-of-law rules.
\end{itemize}
}

\todoitem{Formallize Happy Court Rule}

\doctrine{Doctrinal steps for happy court rule AKA \textit{Feren's} Rule}
{
\todoitem{Clean this up}
\begin{itemize}
    \item Under the Happy Court rule, court should pretend like they are transferor court
    \item Will not use their own choice of law rules (from common law)
    \item Use transferor law: characterize procedural or substantive
    \item if procedural, use transferor law. 
\end{itemize}
}

\subsection{If the federal court does not need to transfer, but may transfer, should it transfer?}

\casebox{\textit{Atlantic Marine Construction Co. v. U.S. District Court}}{571 U.S. 49 (2013)}
{
\textbf{Takeaway}
\begin{itemize}
    \item A forum-selection clause does not make venue improper under \S\,1391 or Rule 12(b)(3).
    \item A clause selecting a federal forum is enforced through the transfer statutes. A clause selecting a state or foreign forum is enforced through forum non conveniens.
    \item When a valid forum-selection clause applies:
    \begin{itemize}
        \item the plaintiff's choice of the current forum receives no deference;
        \item private-interest factors are treated as favoring the selected forum; and
        \item only public-interest factors may ordinarily defeat transfer.
    \end{itemize}
    \item Transfer ordinarily should occur unless extraordinary public-interest circumstances strongly disfavor the selected forum.
    \item The transferee court does not automatically apply the transferor court's choice-of-law rules because a plaintiff who violated the clause has no legitimate venue privilege to preserve.
\end{itemize}
}

\chapter{Erie Doctrine}

\statutebox{28 U.S.C. \S\,1652}
{
The Rules of Decision Act provides that, except where the Constitution, treaties, or federal statutes require otherwise, the laws of the several states shall be treated as rules of decision in federal civil actions where they apply. Under \textit{Erie}, this includes state statutory and state common law. The statute does not authorize a general federal common law governing state-law issues.
}

\doctrine{Holder in Due Course Doctrine}
{
A holder in due course is a person who takes a negotiable instrument for value, in good faith, and without notice of defects or defenses. The holder may generally enforce the instrument against the party who issued or accepted it, even when the underlying transaction was defective.

In \textit{Swift v. Tyson}, the Court treated this as a question of general commercial law rather than state law. \textit{Erie} later rejected the idea that federal courts could create and apply a general federal common law in diversity cases.
}

\casebox{\textit{Swift v. Tyson}}{41 U.S. (16 Pet.) 1 (1842)}
{
\textbf{Takeaway}
\begin{itemize}
    \item Before \textit{Erie}, federal courts sitting in diversity applied federal ``general law'' to questions of general commercial law.
    \item The Court interpreted the Rules of Decision Act as requiring federal courts to follow state statutes and local rules, but not state judicial decisions concerning general commercial law.
    \item The Court applied the holder-in-due-course rule, under which a person who takes a negotiable instrument for value, in good faith, and without notice of defects may enforce it against the other party.
    \item \textit{Swift} created the federal general-common-law regime, which allowed parties to obtain different substantive results by choosing federal court.
    \item \textit{Erie Railroad Co. v. Tompkins} later rejected this approach.
\end{itemize}
}

\casebox{\textit{Erie Railroad Co. v. Tompkins}}{304 U.S. 64 (1938)}
{
\textbf{Takeaway}
\begin{itemize}
    \item \textit{Erie} overruled \textit{Swift} and rejected the existence of federal general common law.
    \item In diversity cases, federal courts apply state substantive law unless the Constitution or Congress provides a federal rule.
    \item Federal courts generally apply federal procedural law and state substantive law.
    \item The decision was motivated by federalism, the need to prevent forum shopping, and the need to prevent different substantive results based solely on whether a party selected state or federal court.
\end{itemize}
}

\doctrine{Erie Doctrine}
{
In diversity cases, federal courts apply state substantive law and federal procedural law. When no federal constitutional provision, statute, or procedural rule controls, the court applies the law that a state court in the same state would apply.

The doctrine is designed to prevent forum shopping and inequitable administration of the law. It also respects the authority of states to define their own substantive rights and obligations.
}


\end{document}
